






% % \vspace{-3mm}
% \subsubsection{Planning Effects}
% \label{subsec:planning_results}



% We evaluate the effect of planning on reasoning accuracy, robustness, and runtime using the benchmark suite described in Section~\ref{subsec:bench_planning}. In contemporary agent frameworks, planning is implemented as an LLM-mediated process in which the model first generates an explicit plan that is then injected into the context to guide answer generation. To isolate the impact of planning interface design, we compare three execution modes under identical models and decoding settings: \emph{NoPlan}, where the model answers directly; \emph{Crew-Plan}, implemented via CrewAI, which enforces a rigid, schema-constrained two-stage planning interface; and \emph{Direct-LLM-Plan}, which allows free-form plan generation without syntactic constraints. For MATH, we evaluate a 100-problem subset preserving the original complexity distribution. Observed differences therefore isolate planning orchestration effects from the LLM model.



% Table~\ref{tab:reasoning-accuracy} shows that schema-constrained planning consistently degrades accuracy across GSM8K, CSQA, and MATH-100 for both local and remote models, \emph{across evaluated benchmarks}. For example, under Crew-Plan, Qwen-7B drops from \emph{13.0\%} to \emph{3.4\%} on GSM8K and from \emph{69.9\%} to \emph{37.6\%} on CSQA, while GPT-OSS-20B declines from \emph{94.6\%} to \emph{91.6\%} on GSM8K and from \emph{80.0\%} to \emph{48.0\%} on MATH-100. Similar degradation appears across both local and remote models. In contrast, Direct-LLM-Plan preserves or improves accuracy in most settings, with Qwen-7B improving to \emph{28.7\%} on GSM8K and Llama-3.1-8B increasing from \emph{65.6\%} to \emph{71.9\%}, indicating that planning is beneficial when not syntactically constrained.


% The underlying cause is revealed by Table~\ref{tab:reasoning-failures} and Table~\ref{tab:reasoning-time-mult}. Crew-Plan exhibits formatting failure rates, with Qwen-7B failing on \emph{84.7\%} of GSM8K instances and \emph{70.0\%} of MATH-100 cases. These failures arise from schema deviations that prevent parsing or execution of otherwise valid reasoning traces. Failures dominate accuracy loss, not reasoning errors, across evaluated datasets. Schema enforcement also introduces substantial overhead: Crew-Plan increases runtime by \emph{7.4×} on GSM8K and \emph{18.5×} on CSQA for GPT-OSS-20B, exceeds \emph{30×} for Llama-3.1-8B on CSQA, and incurs \emph{2×–7×} slowdowns even for large remote models. By contrast, Direct-LLM-Plan yields a \emph{0.0\%} failure rate across all models, because no output format is required, and adds modest overhead (\emph{1.2×–6.6×}), reflecting the cost of generating an additional reasoning trace without orchestration bottlenecks. Overall, planning effectiveness is determined by interface design rather than model capability alone in deployments: free-form planning preserves robustness with predictable cost, while schema-constrained planning degrades both accuracy and efficiency.




% \input{tables/Experiement-Single-Agent-Evaluation-Planning}






\vspace{-1mm}
\subsubsection{Planning Effects}
\label{subsec:planning_results}
We evaluate how the \emph{planning interface} in a framework affects reasoning accuracy, robustness, and runtime, separating architectural constraints from the LLM’s ability to generate plans. Using the MAFBench planning pipeline (Section~\ref{subsec:bench_planning}), we fix the benchmarks (GSM8K~\cite{cobbe2021training}, CSQA~\cite{talmor-etal-2019-commonsenseqa}, MATH~\cite{hendrycks2measuring}), prompts, decoding settings, and scoring. We verify that observed effects are architectural rather than model-specific by running all interfaces across multiple LLM backends. We vary only the framework-level planning interface: \emph{NoPlan} (direct answer), \emph{Crew-Plan} (schema-constrained two-stage planning in CrewAI~\cite{crewai2025}), and \emph{Direct-LLM-Plan} (free-form plan text injected before execution). For MATH, we evaluate a 100-problem subset preserving the original complexity distribution. 

Table~\ref{tab:reasoning-accuracy} shows that schema-constrained planning reduces accuracy across datasets and models, while free-form planning preserves or improves accuracy in most settings. For example, under Crew-Plan, Qwen-7B drops from \emph{13.0\%} to \emph{3.4\%} on GSM8K and from \emph{69.9\%} to \emph{37.6\%} on CSQA, and GPT-OSS-20B drops from \emph{80.0\%} to \emph{48.0\%} on MATH-100. In contrast, Direct-LLM-Plan improves Qwen-7B to \emph{28.7\%} on GSM8K and improves Llama-3.1-8B from \emph{65.6\%} to \emph{71.9\%}, showing that adding a planning stage can help \emph{in practice} when the interface does not impose brittle structure.


Table~\ref{tab:reasoning-failures} and Table~\ref{tab:reasoning-time-mult} identify the architectural cause. Crew-Plan introduces large formatting failure rates, where the LLM fails to follow the framework’s required output schema in practice. When the schema is violated, the framework cannot parse or execute the generated plan, even if the underlying reasoning is correct, leading to \emph{84.7\%} failures on GSM8K and \emph{70.0\%} on MATH-100 for Qwen-7B, which dominate the observed accuracy loss. Crew-Plan also imposes high orchestration overhead by adding an extra LLM call together with schema generation, validation, and parsing, increasing runtime by \emph{7.4$\times$} on GSM8K and \emph{18.5$\times$} on CSQA for GPT-OSS-20B, and exceeding \emph{30$\times$} for Llama-3.1-8B on CSQA. In contrast, Direct-LLM-Plan yields \emph{0.0\%} formatting failures across models. It adds only the additional cost of generating a plan, producing smaller slowdowns (\emph{1.2$\times$--6.6$\times$}) and avoiding any parsing bottleneck.

% These results show that planning outcomes in deployed agent systems are governed primarily by \emph{interface design} rather than by LLM planning ability alone. When building a multi-agent system, planning should be implemented as a permissive stage that tolerates variability in plan text and avoids hard parsing constraints, otherwise the framework converts valid reasoning into failures, reduces robustness, and introduces large runtime overhead that limits scalability.

These results show that planning outcomes in deployed agent systems are driven primarily by \emph{interface design}, not LLM planning ability alone. Planning should therefore be implemented as a permissive stage that tolerates variability in plan text; rigid parsing converts valid reasoning into failures, reduces robustness, and introduces substantial runtime overhead that limits scalability.




% \setlength{\tabcolsep}{8pt}
% \renewcommand{\arraystretch}{1}
% \begin{table*}[t]
% \centering
% \footnotesize
% \caption{Reasoning results across benchmarks with different LLMs. 
% This experiment is designed to measure the effectiveness of CrewAI’s \emph{reasoning} feature by comparing performance in two modes: 
% \xmark = execution without reasoning, 
% \checkmark = with predefined reasoning, 
% and F = cases where the model failed to follow CrewAI’s required reasoning output format. 
% Metrics reported include accuracy, runtime (mean / total), and token usage.}% \textcolor{red}{check https://arxiv.org/pdf/2502.06772}}
% \vspace{-3mm}
% \label{tab:reasoning-results}
% \begin{tabular}{llcc|cc|cc}
% \toprule
% \multirow{2}{*}{\textbf{LLM}} & \multirow{2}{*}{\textbf{Metric}} 
% & \multicolumn{2}{c|}{\textbf{GSM8K}} 
% & \multicolumn{2}{c|}{\textbf{CSQA}} 
% & \multicolumn{2}{c}{\textbf{MATH-100}} \\
% \cmidrule(lr){3-4} \cmidrule(lr){5-6} \cmidrule(lr){7-8}
%  & & \xmark & \checkmark & \xmark & \checkmark & \xmark & \checkmark \\
% \midrule
% \multicolumn{8}{c}{\textbf{Local Models}} \\
% \midrule

% \multirow{3}{*}{DeepSeek-LLM:7B} 
%  & Runtime (sec) & 
%      \minmaxmean{0.43}{478.63}{18.98}{25038.7} & \textbf{F} & % GSM8K
%      \minmaxmean{0.32}{44.59}{3.24}{3951.8} & \textbf{F} & % CSQA
%      \minmaxmeanTODO{0}{100}{20.00}{10000} & \textbf{F} \\ % MATH

%  & Tokens        & 
%      \minmaxmean{0}{5720}{128.24}{169145} & \textbf{F} & % GSM8K
%      \minmaxmean{1}{661}{23.80}{29055} & \textbf{F} & % CSQA
%      \minmaxmeanTODO{0}{100}{20.00}{10000} & \textbf{F} \\ % MATH
 
%  & Accuracy (\%) & 
%     23.35 & \textbf{F} & % GSM8K
%     40.95 & \textbf{F} & % CSQA
%     \textcolor{red}{X} & \textbf{F} \\ % MATH

% \midrule

% \multirow{3}{*}{Llama3.1:7B} 
%  & Runtime (sec) & 
%      \minmaxmean{0.67}{600.26}{17.82}{23503.30} & \minmaxmean{4.75}{908.76}{74.11}{97753.05} & % GSM8K
%      \minmaxmean{1.45}{2.69}{1.58}{1925.22} & \minmaxmean{2.45}{80.93}{12.98}{15848.58} & % CSQA
%      \minmaxmean{0.46}{103.30}{27.57}{2757} & \minmaxmeanTODO{0}{100}{20.00}{10000} \\ % MATH

%  & Tokens        & 
%      \minmaxmean{0}{4220}{121.55}{160324} & \minmaxmean{0}{17421}{92.67}{122230} & % GSM8K
%      \minmaxmean{1}{36}{1.27}{1554} & \minmaxmean{0}{81}{1.43}{1740} & % CSQA
%      \minmaxmean{1}{1087}{224.59}{22459} & \minmaxmeanTODO{0}{100}{20.00}{10000} \\ % MATH
 
%  & Accuracy (\%) & 
%      52.84 & \accChange{52.84}{50.64} $^{(8\%F)}$ & % GSM8K
%      66.34 & \accChange{66.34}{66.50} & % CSQA
%      0.07 & \textcolor{red}{X} \\ % MATH

% \midrule

% \multirow{3}{*}{Qwen:7b} 
%  & Runtime (sec) & 
%      \minmaxmean{0.56}{32}{5.31}{7000} & \minmaxmean{8.11}{173.15}{28.82}{38019} & % GSM8K
%      \minmaxmean{0.46}{10.97}{1.13}{1381} & \minmaxmean{7.9}{189.9}{28.90}{35287} & % CSQA
%      \minmaxmeanTODO{0}{100}{20.00}{10000} & \minmaxmeanTODO{0}{100}{20.00}{10000} \\ % MATH

%  & Tokens        & 
%      \minmaxmean{1}{738}{82.75}{109159} & \minmaxmean{0}{427}{43.96}{57989} & % GSM8K
%      \minmaxmean{1}{243}{9.20}{11238} & \minmaxmean{0}{253}{11.84}{14457} & % CSQA
%      \minmaxmeanTODO{0}{100}{20.00}{10000} & \minmaxmeanTODO{0}{100}{20.00}{10000} \\ % MATH
 
%  & Accuracy (\%) & 
%      11.98 & \accChange{11.98}{5.53} $^{(56\%F)}$ & % GSM8K
%      25.39 & \accChange{25.39}{18.42} & % CSQA
%      \textcolor{red}{X} & \textcolor{red}{X} \\ % MATH
     
% \midrule

% \multicolumn{8}{c}{\textbf{Remote Models}} \\
% \midrule

% \multirow{4}{*}{GPT-o4-Mini} 
%  & Runtime (sec) & 
%  %\minmaxmean{minValue}{maxValue}{meanValue}{totalValue}
%      \minmaxmean{0.44}{20.85}{3.17}{4180.6} & \minmaxmean{3.18}{68.02}{7.02}{9253.6} & % GSM8K
%      \minmaxmean{0.40}{35.91}{1.18}{1444.8} & \minmaxmean{3.17}{16.57}{5.58}{6818.8} & % CSQA
%      \minmaxmean{0.64}{41.05}{12.09}{1209} & \minmaxmean{5.33}{469.90}{25.12}{2512.7} \\ % MATH

%  & Tokens        & 
%      \minmaxmean{1}{7}{1.01}{1330} & \minmaxmean{1}{9}{1.02}{1340} & % GSM8K
%      \minmaxmean{1}{3}{1.00}{1225} & \minmaxmean{1}{1}{1.00}{1221} & % CSQA
%      \minmaxmean{1}{612}{11.57}{1157} & \minmaxmeanTODO{0}{100}{20.00}{220} \\ % MATH

%  & Accuracy (\%) & 
%      81.50 & \accChange{81.50}{87.57} & % GSM8K
%      82.23 & \accChange{82.23}{81.65} & % CSQA
%      60.00 & \accChange{60.00}{63.00} \\ % MATH

% \midrule

% \multirow{3}{*}{GPT-4.1} 
%  & Runtime (sec) & 
%      \minmaxmean{0.59}{9.06}{1.99}{2625.50} & 
%      \minmaxmean{3.07}{37.86}{6.85}{9036.44} & % GSM8K
%      \minmaxmean{0.55}{13.75}{1.94}{2371.87} & 
%      \minmaxmean{3.30}{26.72}{6.51}{7950.58} & % CSQA
%      \minmaxmeanTODO{0}{100}{20.00}{10000} & 
%      \minmaxmeanTODO{0}{100}{20.00}{10000} \\ % MATH

%  & Tokens        & 
%      \minmaxmean{1}{2}{1}{1322} & \minmaxmean{1}{2}{1}{1320} & % GSM8K
%      \minmaxmean{1}{1}{1}{1221} & \minmaxmean{1}{1}{1}{1221} & % CSQA
%      \minmaxmeanTODO{0}{100}{20.00}{10000} & 
%      \minmaxmeanTODO{0}{100}{20.00}{10000} \\ % MATH
 
%  & Accuracy (\%) & 
%      90.90 & \accChange{90.90}{78.54} & % GSM8K
%      87.06 & \accChange{87.06}{87.31} & % CSQA
%      \textcolor{red}{X} & \textcolor{red}{X} \\ % MATH

     
% \bottomrule
% \end{tabular}
% \end{table*}

% \setlength{\tabcolsep}{4pt}
% \begin{table*}[t]
% \centering
% \footnotesize
% \caption{Reasoning results across CrewAI (no-planning vs planning) and Direct-LLM planning.}
% \vspace{-3mm}
% \label{tab:reasoning-results}
% \begin{tabular}{llccc|ccc|ccc}
% \toprule
% \multirow{2}{*}{\textbf{LLM}} & \multirow{2}{*}{\textbf{Metric}} & \multicolumn{3}{c|}{\textbf{GSM8K}} & \multicolumn{3}{c|}{\textbf{CSQA}} & \multicolumn{3}{c}{\textbf{MATH-100}} \\
% \cmidrule(lr){3-5} \cmidrule(lr){6-8} \cmidrule(lr){9-11}
%  & & NoPlan & CrewPlan & Direct & NoPlan & CrewPlan & Direct & NoPlan & CrewPlan & Direct \\
% \midrule
% \multicolumn{11}{c}{\textbf{Local Models}} \\
% \midrule
% \multirow{3}{*}{deepseek-llm:7b} 
%  & Runtime (sec)  & \minmaxmean{0.75}{164.78}{18.15}{23935.76} & \minmaxmean{8.56}{647.44}{56.43}{74437.02} & \minmaxmean{9.75}{254.93}{47.08}{62101.65} & \minmaxmean{0.33}{52.19}{3.31}{4037.68} & \minmaxmean{5.34}{314.34}{43.96}{53678.87} & \minmaxmean{2.35}{105.27}{21.87}{26706.73} & \minmaxmean{2.44}{206.39}{54.10}{5410.02} & \minmaxmean{37.35}{983.63}{153.69}{15369.09} & \minmaxmean{8.26}{223.34}{65.40}{6540.22} \\
%  & Tokens  & \minmaxmean{0.00}{2143.00}{113.25}{149380.00} & \minmaxmean{0.00}{1040.00}{72.88}{96124.00} & \minmaxmean{48.00}{1146.00}{260.74}{343912.00} & \minmaxmean{1.00}{526.00}{22.72}{27736.00} & \minmaxmean{0.00}{698.00}{20.19}{24656.00} & \minmaxmean{1.00}{704.00}{140.67}{171760.00} & \minmaxmean{1.00}{1071.00}{218.23}{21823.00} & \minmaxmean{0.00}{559.00}{97.41}{9741.00} & \minmaxmean{7.00}{876.00}{284.66}{28466.00} \\
%  & Accuracy (\%)  & 23.12 & 15.39$^{(25.47\%F)}$ & 19.64$^{(0.00\%F)}$ & 42.83 & 27.93$^{(0.00\%F)}$ & 35.87$^{(0.00\%F)}$ & 9.00 & 3.00$^{(40.00\%F)}$ & 4.00$^{(0.00\%F)}$ \\
% \midrule
% \multirow{3}{*}{llama3.1:8b} 
%  & Runtime (sec)  & \minmaxmean{3.47}{600.31}{31.96}{42150.01} & \minmaxmean{26.62}{1150.28}{201.13}{265288.72} & \textcolor{red}{X} & \minmaxmean{3.21}{30.15}{3.58}{4369.17} & \minmaxmean{20.82}{814.96}{112.16}{136945.64} & \textcolor{red}{X} & \minmaxmean{3.70}{600.27}{76.38}{7637.65} & \minmaxmean{40.41}{1728.45}{337.49}{33748.87} & \textcolor{red}{X} \\
%  & Tokens  & \minmaxmean{0.00}{1218.00}{110.82}{146174.00} & \minmaxmean{0.00}{2745.00}{65.25}{86062.00} & \textcolor{red}{X} & \minmaxmean{1.00}{50.00}{1.29}{1573.00} & \minmaxmean{0.00}{124.00}{1.62}{1983.00} & \textcolor{red}{X} & \minmaxmean{0.00}{2114.00}{215.25}{21525.00} & \minmaxmean{0.00}{1403.00}{124.40}{12440.00} & \textcolor{red}{X} \\
%  & Accuracy (\%)  & 52.92 & 50.27$^{(10.54\%F)}$ & \textcolor{red}{X}$^{(0\%F)}$ & 68.14 & 63.72$^{(0.00\%F)}$ & \textcolor{red}{X}$^{(0\%F)}$ & 25.00 & 24.00$^{(19.00\%F)}$ & \textcolor{red}{X}$^{(0\%F)}$ \\
% \midrule
% \multirow{3}{*}{phi4:14b} 
%  & Runtime (sec)  & \minmaxmean{20.23}{173.06}{59.13}{77988.88} & \minmaxmean{104.47}{600.23}{208.62}{275165.92} & \textcolor{red}{X} & \minmaxmean{22.50}{90.03}{46.68}{56998.23} & \minmaxmean{115.67}{600.24}{190.98}{233189.35} & \textcolor{red}{X} & \minmaxmean{30.70}{345.57}{110.75}{11074.69} & \minmaxmean{121.78}{666.89}{327.62}{32762.45} & \textcolor{red}{X} \\
%  & Tokens  & \minmaxmean{1.00}{98.00}{1.09}{1442.00} & \minmaxmean{0.00}{726.00}{1.81}{2390.00} & \textcolor{red}{X} & \minmaxmean{1.00}{128.00}{1.40}{1714.00} & \minmaxmean{0.00}{528.00}{3.43}{4187.00} & \textcolor{red}{X} & \minmaxmean{1.00}{1079.00}{21.09}{2109.00} & \minmaxmean{0.00}{605.00}{14.55}{1455.00} & \textcolor{red}{X} \\
%  & Accuracy (\%)  & 89.46 & 82.49$^{(6.14\%F)}$ & \textcolor{red}{X}$^{(0\%F)}$ & 75.68 & 81.33$^{(0.00\%F)}$ & \textcolor{red}{X}$^{(0\%F)}$ & 72.00 & 50.00$^{(21.00\%F)}$ & \textcolor{red}{X}$^{(0\%F)}$ \\
% \midrule
% \multirow{3}{*}{qwen:7b} 
%  & Runtime (sec)  & \minmaxmean{2.63}{116.26}{20.43}{26943.80} & \minmaxmean{0.00}{711.41}{45.79}{60399.04} & \textcolor{red}{X} & \minmaxmean{2.55}{46.37}{5.94}{7253.45} & \minmaxmean{33.07}{475.78}{107.90}{131743.05} & \textcolor{red}{X} & \minmaxmean{9.06}{263.54}{50.11}{5010.91} & \minmaxmean{33.60}{543.87}{146.13}{14613.38} & \textcolor{red}{X} \\
%  & Tokens  & \minmaxmean{1.00}{519.00}{80.32}{105947.00} & \minmaxmean{0.00}{926.00}{16.16}{21320.00} & \textcolor{red}{X} & \minmaxmean{1.00}{280.00}{9.44}{11526.00} & \minmaxmean{0.00}{281.00}{11.42}{13944.00} & \textcolor{red}{X} & \minmaxmean{1.00}{1161.00}{182.41}{18241.00} & \minmaxmean{0.00}{720.00}{65.05}{6505.00} & \textcolor{red}{X} \\
%  & Accuracy (\%)  & 10.99 & 2.88$^{(84.69\%F)}$ & \textcolor{red}{X}$^{(0\%F)}$ & 25.88 & 15.72$^{(0.00\%F)}$ & \textcolor{red}{X}$^{(0\%F)}$ & 9.00 & 4.00$^{(70.00\%F)}$ & \textcolor{red}{X}$^{(0\%F)}$ \\
% \midrule
% \multicolumn{11}{c}{\textbf{Remote Models}} \\
% \midrule
% \multirow{3}{*}{Groq:gpt:oss:20b} 
%  & Runtime (sec)  & \minmaxmean{0.17}{3.90}{0.46}{602.71} & \minmaxmean{0.57}{55.35}{3.38}{4457.47} & \textcolor{red}{X} & \minmaxmean{0.17}{5.05}{0.40}{494.24} & \minmaxmean{0.58}{55.12}{7.49}{9149.41} & \textcolor{red}{X} & \minmaxmean{0.21}{5.29}{1.22}{121.94} & \minmaxmean{0.57}{61.75}{13.94}{1394.22} & \textcolor{red}{X} \\
%  & Tokens  & \minmaxmean{0.00}{2.00}{1.00}{1316.00} & \minmaxmean{0.00}{4.00}{0.98}{1290.00} & \textcolor{red}{X} & \minmaxmean{0.00}{1.00}{0.98}{1202.00} & \minmaxmean{0.00}{1.00}{0.99}{1205.00} & \textcolor{red}{X} & \minmaxmean{0.00}{15.00}{1.29}{129.00} & \minmaxmean{0.00}{17.00}{0.72}{72.00} & \textcolor{red}{X} \\
%  & Accuracy (\%)  & 94.09 & 89.23$^{(2.43\%F)}$ & \textcolor{red}{X}$^{(0\%F)}$ & 81.41 & 81.24$^{(0.00\%F)}$ & \textcolor{red}{X}$^{(0\%F)}$ & 80.00 & 48.00$^{(49.00\%F)}$ & \textcolor{red}{X}$^{(0\%F)}$ \\
% \midrule
% \multirow{3}{*}{gpt-4.1} 
%  & Runtime (sec)  & \minmaxmean{0.90}{27.36}{2.74}{3612.95} & \minmaxmean{2.98}{84.95}{6.60}{8702.99} & \textcolor{red}{X} & \minmaxmean{0.71}{41.95}{2.63}{3208.67} & \minmaxmean{3.61}{52.18}{7.64}{9325.96} & \textcolor{red}{X} & \minmaxmean{0.77}{120.84}{11.44}{1144.02} & \minmaxmean{4.60}{125.64}{13.75}{1375.10} & \textcolor{red}{X} \\
%  & Tokens  & \minmaxmean{1.00}{23.00}{1.03}{1357.00} & \minmaxmean{1.00}{53.00}{1.04}{1377.00} & \textcolor{red}{X} & \minmaxmean{1.00}{1.00}{1.00}{1221.00} & \minmaxmean{1.00}{1.00}{1.00}{1221.00} & \textcolor{red}{X} & \minmaxmean{1.00}{70.00}{3.75}{375.00} & \minmaxmean{1.00}{26.00}{3.02}{302.00} & \textcolor{red}{X} \\
%  & Accuracy (\%)  & 90.67 & 78.70$^{(0.00\%F)}$ & \textcolor{red}{X}$^{(0\%F)}$ & 87.06 & 87.31$^{(0.00\%F)}$ & \textcolor{red}{X}$^{(0\%F)}$ & 86.00 & 60.00$^{(0.00\%F)}$ & \textcolor{red}{X}$^{(0\%F)}$ \\
% \midrule
% \multirow{3}{*}{gpt-4o-mini} 
%  & Runtime (sec)  & \minmaxmean{0.54}{49.53}{5.61}{7402.36} & \minmaxmean{2.54}{45.27}{8.37}{11040.59} & \textcolor{red}{X} & \minmaxmean{0.49}{7.68}{1.57}{1922.32} & \minmaxmean{2.34}{27.02}{7.52}{9185.59} & \textcolor{red}{X} & \minmaxmean{0.66}{35.55}{11.63}{1162.81} & \minmaxmean{6.72}{66.62}{21.18}{2118.43} & \textcolor{red}{X} \\
%  & Tokens  & \minmaxmean{1.00}{197.00}{1.15}{1516.00} & \minmaxmean{1.00}{6.00}{1.01}{1326.00} & \textcolor{red}{X} & \minmaxmean{1.00}{5.00}{1.00}{1227.00} & \minmaxmean{1.00}{1.00}{1.00}{1221.00} & \textcolor{red}{X} & \minmaxmean{1.00}{19.00}{2.37}{237.00} & \minmaxmean{1.00}{330.00}{6.08}{608.00} & \textcolor{red}{X} \\
%  & Accuracy (\%)  & 81.80 & 86.73$^{(0.00\%F)}$ & \textcolor{red}{X}$^{(0\%F)}$ & 81.65 & 81.33$^{(0.00\%F)}$ & \textcolor{red}{X}$^{(0\%F)}$ & 60.00 & 53.00$^{(0.00\%F)}$ & \textcolor{red}{X}$^{(0\%F)}$ \\
% \midrule
% \bottomrule
% \end{tabular}
% \end{table*}








\setlength{\tabcolsep}{1.7pt}
\begin{table}[t]
\centering
\scriptsize
\caption{Accuracy (\%) under three planning modes:
\xmark: No plan,\; \checkmark(Crew): schema-constrained plan,\; \checkmark(LLM): free-form plan.}
\vspace{-3mm}
\renewcommand{\arraystretch}{0.65}
\label{tab:reasoning-accuracy}
\begin{tabular}{l|ccc|ccc|ccc}
\toprule
\multirow{2}{*}{\textbf{LLM}} & \multicolumn{3}{c}{\textbf{GSM8K}~\cite{cobbe2021training}} & \multicolumn{3}{c}{\textbf{CSQA}~\cite{talmor-etal-2019-commonsenseqa}} & \multicolumn{3}{c}{\textbf{MATH-100}~\cite{hendrycks2measuring}} \\
\cmidrule(lr){2-4} \cmidrule(lr){5-7} \cmidrule(lr){8-10}
 & \xmark & \checkmark (crew) & \checkmark (LLM) & \xmark & \checkmark (crew) & \checkmark (LLM)& \xmark & \checkmark (crew) & \checkmark (LLM) \\
\midrule
\rowcolor{gray!20} \multicolumn{10}{c}{\textbf{Local Models}} \\
GPT-OSS-20B~\cite{openai2025gptoss} & 94.6 & \accChange{94.6}{91.6} & \accChange{94.6}{87.6} & 81.4 & \accChange{81.4}{81.2} & \accChange{81.4}{82.1} & 80.0 & \accChange{80.0}{48.0} & \accChange{80.0}{76.0} \\
Phi-4-14B~\cite{abdin2024phi4} & 93.3 & \accChange{93.3}{86.7} & \accChange{93.3}{83.5} & 81.7 & \accChange{81.7}{82.0} & \accChange{81.7}{74.6} & 72.0 & \accChange{72.0}{50.0} & \accChange{72.0}{69.0} \\
Llama-3.1-8B~\cite{dubey2024llama3} & 65.6 & \accChange{65.6}{60.3} & \accChange{65.6}{71.9} & 70.1 & \accChange{70.1}{65.4} & \accChange{70.1}{65.4} & 25.0 & \accChange{25.0}{24.0} & \accChange{25.0}{37.0} \\
Qwen-7B~\cite{bai2024qwenvl} & 13.0 & \accChange{13.0}{3.4} & \accChange{13.0}{28.7} & 69.9 & \accChange{69.9}{37.6} & \accChange{69.9}{74.0} & 9.0 & \accChange{9.0}{4.0} & \accChange{9.0}{13.0} \\

DeepSeek-7B~\cite{deepseek2024v3} & 27.0 & \accChange{27.0}{18.0} & \accChange{27.0}{23.5} & 46.3 & \accChange{46.3}{31.2} & \accChange{46.3}{46.1} & 9.0 & \accChange{9.0}{3.0} & \accChange{9.0}{4.0} \\
\rowcolor{gray!20} \multicolumn{10}{c}{\textbf{Remote Models}} \\
GPT-4.1~\cite{openai2024gpt4} & 94.2 & \accChange{94.2}{80.7} & \accChange{94.2}{78.2} & 87.1 & \accChange{87.1}{87.3} & \accChange{87.1}{87.0} & 86.0 & \accChange{86.0}{60.0} & \accChange{86.0}{84.0} \\
GPT-4o-Mini~\cite{openai2024gpt4} & 86.5 & \accChange{86.5}{90.8} & \accChange{86.5}{84.2} & 81.8 & \accChange{81.8}{81.3} & \accChange{81.8}{83.1} & 60.0 & \accChange{60.0}{53.0} & \accChange{60.0}{59.0} \\
\bottomrule
\end{tabular}
\vspace{-2mm}
\end{table}







\setlength{\tabcolsep}{3.5pt}
\begin{table}[t]
\centering
\scriptsize
% \vspace{-3mm}
\caption{Formatting failures under schema-constrained and free-form planning, isolating interface-induced errors.}
\vspace{-3mm}
\renewcommand{\arraystretch}{0.95}
\label{tab:reasoning-failures}
\begin{tabular}{l|cc|cc|cc}
\toprule
\multirow{2}{*}{\textbf{LLM}} & \multicolumn{2}{c}{\textbf{GSM8K}} & \multicolumn{2}{c}{\textbf{CSQA}} & \multicolumn{2}{c}{\textbf{MATH-100}} \\
\cmidrule(lr){2-3} \cmidrule(lr){4-5} \cmidrule(lr){6-7}
 & \checkmark (CrewAI) & \checkmark (LLM) & \checkmark (CrewAI) & \checkmark (LLM) & \checkmark (CrewAI) & \checkmark (LLM) \\
\midrule
\rowcolor{gray!20} \multicolumn{7}{c}{\textbf{Local Models}} \\
GPT-OSS-20B & 2.4\% & 0.0\% & 1.3\% & 0.0\% & 49.0\% & 13.0\% \\
Phi-4-14B & 6.1\% & 0.0\% & 0.5\% & 0.0\% & 21.0\% & 0.0\% \\
Llama-3.1-8B & 10.5\% & 0.0\% & 4.8\% & 0.0\% & 19.0\% & 0.0\% \\
Qwen-7B & 84.7\% & 0.0\% & 41.0\% & 0.0\% & 70.0\% & 0.0\% \\
DeepSeek-7B & 25.5\% & 0.0\% & 31.0\% & 0.0\% & 40.0\% & 0.0\% \\
\rowcolor{gray!20} \multicolumn{7}{c}{\textbf{Remote Models}} \\
GPT-4.1 & 0.0\% & 0.0\% & 0.0\% & 0.0\% & 0.0\% & 0.0\% \\
GPT-4o-Mini & 0.0\% & 0.0\% & 0.0\% & 0.0\% & 0.0\% & 0.0\% \\
\bottomrule
\end{tabular}
\end{table}




\setlength{\tabcolsep}{5.pt}
\begin{table}[t]
\centering
\scriptsize
% \vspace{-4mm}
\caption{Runtime multiplier (relative to NoPlan) under schema-constrained and free-form planning. Lower is faster.}
\vspace{-3mm}
\renewcommand{\arraystretch}{0.9}
\label{tab:reasoning-time-mult}
\begin{tabular}{l|cc|cc|cc}
\toprule
\multirow{2}{*}{\textbf{LLM}} & \multicolumn{2}{c}{\textbf{GSM8K}} & \multicolumn{2}{c}{\textbf{CSQA}} & \multicolumn{2}{c}{\textbf{MATH}} \\
\cmidrule(lr){2-3} \cmidrule(lr){4-5} \cmidrule(lr){6-7}
 & $\Delta$ CrewAI & $\Delta$ LLM & $\Delta$ CrewAI & $\Delta$ LLM & $\Delta$ CrewAI & $\Delta$ LLM \\
\midrule
\rowcolor{gray!20} \multicolumn{7}{c}{\textbf{Local Models}} \\
GPT-OSS-20B & $\times$ 7.4 & $\times$ 2.3 & $\times$ 18.51 & $\times$ 1.7 & $\times$ 11.43 & $\times$ 1.7 \\
Phi-4-14B & $\times$ 3.5 & $\times$ 1.4 & $\times$ 4.1 & $\times$ 1.6 & $\times$ 3.0 & $\times$ 1.7 \\
Llama-3.1-8B & $\times$ 6.3 & $\times$ 1.3 & $\times$ 31.34 & $\times$ 4.2 & $\times$ 4.4 & $\times$ 1.3 \\
Qwen-7B & $\times$ 2.2 & $\times$ 2.5 & $\times$ 18.16 & $\times$ 0.9 & $\times$ 2.9 & $\times$ 1.8 \\
DeepSeek-7B & $\times$ 3.1 & $\times$ 2.6 & $\times$ 13.29 & $\times$ 6.6 & $\times$ 2.8 & $\times$ 1.2 \\
\rowcolor{gray!20} \multicolumn{7}{c}{\textbf{Remote Models}} \\
GPT-4.1 & $\times$ 2.4 & $\times$ 1.5 & $\times$ 2.9 & $\times$ 1.1 & $\times$ 1.2 & $\times$ 1.1 \\
GPT-4o-Mini & $\times$ 1.5 & $\times$ 1.4 & $\times$ 4.8 & $\times$ 2.3 & $\times$ 1.8 & $\times$ 1.4 \\
\bottomrule
\end{tabular}
\end{table}